\documentclass[12pt]{article}
\usepackage{amsmath,amssymb}
\usepackage{geometry}
\usepackage{setspace}
\geometry{margin=1in}
\setstretch{1.15}

\begin{document}

\begin{center}
\Large \textbf{Lecture Scribe: Bayes’ Theorem, Random Variables, and Probability Mass Function}

\vspace{0.5em}

\normalsize
\textbf{Course:} CSE 400 – Fundamentals of Probability in Computing \\
\textbf{Lecture:} 5 \\
\textbf{Instructor:} Dhaval Patel, PhD \\
\textbf{Date:} January 20, 2026
\end{center}

\hrule
\vspace{1em}

\section*{1. Introduction}

This lecture covers:
\begin{itemize}
\item Bayes’ Theorem as a weighted average of conditional probabilities
\item Learning Bayes’ Theorem through examples
\item Formal introduction to the Law of Total Probability and Bayes’ Formula
\item Definition and concept of Random Variables
\item Probability Mass Function (PMF) for discrete random variables
\end{itemize}

\hrule
\vspace{1em}

\section*{2. Bayes’ Theorem}

\subsection*{2.1 Weighted Average of Conditional Probabilities}

Let $( A )$ and $( B )$ be events.

We may express:
\[
A = AB \cup AB^c
\]

For an outcome to be in $( A )$, it must either be in both $( A )$ and $( B )$, or be in $( A )$ but not in $( B )$.

Since $( AB )$ and $( AB^c )$ are mutually exclusive, by \textbf{Axiom 3}:
\[
\Pr(A) = \Pr(AB) + \Pr(AB^c)
\]

Using conditional probability:
\[
\Pr(A) = \Pr(A|B)\Pr(B) + \Pr(A|B^c)[1 - \Pr(B)]
\]

\textbf{Result:}  
The probability of event $( A )$ is a \textbf{weighted average} of conditional probabilities, where weights are given by the probabilities of the conditioning events.

\hrule
\vspace{1em}

\subsection*{2.2 Learning by Example}

\subsubsection*{Example 3.1 (Part 1/2)}

An insurance company classifies people into:
\begin{itemize}
\item Accident-prone
\item Not accident-prone
\end{itemize}

Given:
\begin{itemize}
\item $( \Pr(\text{Accident} \mid \text{Accident-prone}) = 0.4 )$
\item $( \Pr(\text{Accident} \mid \text{Not accident-prone}) = 0.2 )$
\item $( \Pr(\text{Accident-prone}) = 0.3 )$
\end{itemize}

Let:
\begin{itemize}
\item $( A_1 )$: Policyholder has an accident within one year
\item $( A )$: Policyholder is accident-prone
\end{itemize}

Then:
\[
\Pr(A_1) = \Pr(A_1|A)\Pr(A) + \Pr(A_1|A^c)\Pr(A^c)
\]
\[
= (0.4)(0.3) + (0.2)(0.7) = 0.26
\]

\subsubsection*{Example 3.1 (Part 2/2)}

Given that a new policyholder \textbf{has had an accident}, find the probability that the person is accident-prone.

\[
\Pr(A|A_1) = \frac{\Pr(AA_1)}{\Pr(A_1)}
\]

Using:
\[
\Pr(AA_1) = \Pr(A)\Pr(A_1|A)
\]

\[
\Pr(A|A_1) = \frac{(0.3)(0.4)}{0.26} = \frac{6}{13}
\]

\hrule
\vspace{1em}

\section*{3. Formal Introduction: Law of Total Probability and Bayes’ Formula}

\subsection*{3.1 Law of Total Probability}

Suppose:
\[
B_1, B_2, \ldots, B_n
\]
are mutually exclusive events such that:
\[
\bigcup_{i=1}^{n} B_i = B
\]

Then:
\[
A = \bigcup_{i=1}^{n} AB_i
\]

Since the events $( AB_i )$ are mutually exclusive:
\[
\Pr(A) = \sum_{i=1}^{n} \Pr(AB_i)
\]

Using:
\[
\Pr(AB_i) = \Pr(A|B_i)\Pr(B_i)
\]

We obtain:
\[
\Pr(A) = \sum_{i=1}^{n} \Pr(A|B_i)\Pr(B_i)
\]

This is known as the \textbf{Law of Total Probability (Formula 3.4)}.

\subsection*{3.2 Bayes’ Formula (Proposition 3.1)}

Using:
\[
\Pr(AB_i) = \Pr(B_i|A)\Pr(A)
\]

We obtain:
\[
\Pr(B_i|A) =
\frac{\Pr(A|B_i)\Pr(B_i)}
{\sum_{j=1}^{n} \Pr(A|B_j)\Pr(B_j)}
\]

This is known as the \textbf{Bayes Formula (Proposition 3.1)}.

Definitions:
\begin{itemize}
\item $( \Pr(B_i) )$: \textbf{a priori probability}
\item $( \Pr(B_i|A) )$: \textbf{a posteriori probability}
\end{itemize}

\subsection*{3.3 Learning by Example}

\subsubsection*{Example 3.2 (Card Problem)}

Three cards:
\begin{itemize}
\item RR (both red)
\item BB (both black)
\item RB (one red, one black)
\end{itemize}

One card is chosen randomly and placed down.  
The \textbf{upper side is red}.

Let:
\begin{itemize}
\item $( RR, BB, RB )$: corresponding card events
\item $( R )$: event that upturned side is red
\end{itemize}

\[
\Pr(RB|R) =
\frac{\Pr(R|RB)\Pr(RB)}
{\Pr(R|RR)\Pr(RR) + \Pr(R|RB)\Pr(RB) + \Pr(R|BB)\Pr(BB)}
\]

\[
\frac{(1/2)(1/3)}
{(1)(1/3) + (1/2)(1/3) + (0)(1/3)}
= \frac{1}{3}
\]

\hrule
\vspace{1em}

\section*{4. Random Variables}

\subsection*{4.1 Definition}

A \textbf{random variable} $( X )$ on a sample space $( \Omega )$ is a function:
\[
X : \Omega \rightarrow \mathbb{R}
\]

It assigns a real number to each sample point $( \omega \in \Omega )$.

Only \textbf{discrete random variables} are considered:
\begin{itemize}
\item Finite or countably infinite range
\item Values form a discrete subset of $( \mathbb{R} )$
\end{itemize}

\subsection*{4.2 Distribution of a Random Variable}

For any value $( a )$ in the range of $( X )$:
\[
\{\omega \in \Omega : X(\omega) = a\}
\]
is an event, denoted by:
\[
X = a
\]

The probability:
\[
\Pr[X = a]
\]
is defined.

The collection of these probabilities for all possible values of $( a )$ is called the \textbf{distribution of $( X )$}.

\subsection*{4.3 Example}

\textbf{Example 1: Tossing 3 Fair Coins}

Let $( Y )$ be the number of heads.

Possible values: $( 0,1,2,3 )$

\[
\Pr(Y=0) = \frac{1}{8}, \quad
\Pr(Y=1) = \frac{3}{8}, \quad
\Pr(Y=2) = \frac{3}{8}, \quad
\Pr(Y=3) = \frac{1}{8}
\]

Since $( Y )$ must take one of these values:
\[
1 = \sum_{i=0}^{3} \Pr(Y=i)
\]

\hrule
\vspace{1em}

\section*{5. Probability Mass Function (PMF)}

\subsection*{5.1 Definition}

A random variable that can take on \textbf{at most a countable number of values} is said to be \textbf{discrete}.

Let $( X )$ be a discrete random variable with range:
\[
R_X = \{x_1, x_2, x_3, \ldots\}
\]

The function:
\[
p(x_k) = \Pr(X = x_k)
\]
is called the \textbf{Probability Mass Function (PMF)} of $( X )$.

Since $( X )$ must take one of the values $( x_k )$:
\[
\sum_k p(x_k) = 1
\]

\subsection*{5.2 Worked Example}

Given:
\[
p(i) = c \frac{\lambda^i}{i!}, \quad i = 0,1,2,\ldots
\]
where $( \lambda > 0 )$.

Since:
\[
\sum_{i=0}^{\infty} p(i) = 1
\]

\[
c \sum_{i=0}^{\infty} \frac{\lambda^i}{i!} = 1
\]

Using:
\[
e^\lambda = \sum_{i=0}^{\infty} \frac{\lambda^i}{i!}
\]

\[
ce^\lambda = 1 \Rightarrow c = e^{-\lambda}
\]

Therefore:
\[
\Pr(X=0) = e^{-\lambda}
\]

\[
\Pr(X>2) = 1 - [\Pr(X=0) + \Pr(X=1) + \Pr(X=2)]
\]

\hrule
\vspace{1em}

\section*{6. Conclusion}

This lecture formally established:
\begin{itemize}
\item Bayes’ Theorem via conditional probabilities
\item Law of Total Probability
\item Bayes’ Formula and its interpretation
\item Definition and distribution of random variables
\item Probability Mass Function for discrete random variables
\end{itemize}

\begin{center}
\textbf{End of Lecture 5 Scribe}
\end{center}

\end{document}
